\documentclass[a4paper,12pt]{article}
\usepackage[T1]{fontenc}
\usepackage[utf8]{inputenc}
\usepackage{t1enc}
\usepackage[dvips]{color}
\usepackage{amsmath,amssymb,makeidx,stmaryrd}
\usepackage[normalem]{ulem}
\usepackage{fancybox}
\definecolor{gris}{gray}{0.9}
\def\bas{\fcolorbox{gris}{gris}{\parbox{0.986\textwidth}{
	\color{red}\footnotesize Mercredi 22 Mars 2000\hfill Page : \thepage/1}}}
\def\haut{\fcolorbox{gris}{gris}{\parbox{0.986\textwidth}{
	\color{red}
	\bf Concours Général de Mathématiques \hfill 2000
	}}}
\makeatletter
\def\ps@cg{%
\let\@oddhead\haut%
\let\@evenhead\haut%
\let\@oddfoot\bas
\let\@evenfoot\bas}
\makeatother
\pagestyle{cg}
\usepackage[left=3.cm, right=3.cm, top=3cm, bottom=3cm]{geometry}
\newcommand{\vect}[1]{\overrightarrow{\,\mathstrut#1\,}}
\parindent0pt
\everymath{\displaystyle}
\usepackage[french]{babel}

\DecimalMathComma
\usepackage[np]{numprint}
\begin{document}

\begin{center}\textbf{{\color{blue} Exercice 1}}
\end{center}

On dispose de $b$ boules blanches et $n$ boules noires - au moins une de
chaque -, que l'on répartit entre deux urnes de façon qu'aucune d'elles ne
soit vide~; on note $s$ le nombre de boules dans la première, et $r$ celui
de ces boules qui sont blanches. L'évènement considéré est le tirage d'une
boule au hasard dans l'une des urnes choisie au hasard~; le but de
l'exercice est de déterminer les répartitions rendant maximale la
probabilité $p$ de tirer une boule blanche.

\medskip

\begin{enumerate}
\item Exprimer $p$ en fonction de $b, n, r$ et $s$.
\item Dans cette question, l'on fixe la valeur de $s$ ;
comment choisir $r$ pour augmenter $p$~?
\item Résoudre l'exercice.
\item Quelles généralisations proposez-vous en augmentant les nombres
de couleurs et d'urnes~?
\end{enumerate}

\begin{center}\textbf{\color{blue} Problème}
\end{center}

Ce problème traite des triangles $ABC$ dits cartésiens, c'est-à-dire à
cotés entiers $BC = a$, $CA = b$ et $AB = c$ dont l'angle en A mesure
$\frac{2\pi}3$ radians. Sauf avis contraire $ABC $ est supposé cartésien.
\begin{enumerate}
\item Notant $H$ son orthocentre orthogonalement projeté en $(U,V,W)$
	sur les trois cotés, déterminer les nombres rationnels parmi
	\begin{center}
$AU,\: BV,\: CW,\: HA,\: HB,\: HC,\: HU, HV,\: HW,\: AW,\: AV,\: BU,\: BW,\: CV $ et $CU$.
	\end{center}
\item Notant $I$ son son centre du cercle inscrit, $J$ l'intersection de
	la bissectrice intérieure en $A$ et des bissectrices extérieures en les
	autres sommets, et $P, Q $ les intersections de la droite $BC$ et des deux
	bissectrices en $A$, déterminer les nombres rationnels parmi 
	
	\begin{center}$PB,\: PC,\: QB,\:
	QC,\: AI,\: AJ,\: AP$ et $AQ$.\end{center}
	
\item On suppose désormais $b$ et $c$ premiers entre eux. Montrer que,
	quitte à échanger $b$ et $c$, $a + b - c$ est multiple de $3$ et
	$a- b + c$  ne l'est pas.
\item On pose $\dfrac {a + b - c}{3c}=\dfrac pq$ où $p$ et $q$ sont des entiers
	strictement positifs premiers entre eux.
	
Notant $d$ le PGCD de $p(3p + 2q)$
	et de $q(2p + q)$, calculer $a, b$ et $c$ en fonction de $p, q$ et $d$.
\item Montrer que $q$ n'est pas multiple de $3$, puis que $d = 1$.
\item En déduire une condition nécéssaire et suffisante pour qu'un
	triangle soit cartésien de cotés premiers entre eux puis, par des remarques
	géométriques, une caractérisation analogue des triangles à cotés entiers
	$BC = a,\ CA = b$ et $AB = c$ premiers entre eux dont l'angle en $A$ mesure
	$\frac \pi 3$ radians.
\end{enumerate}

\begin{center}\textbf{\color{blue} Exercice 2}
\end{center}

Soient $A, B, C$ trois points deux à deux distincts dans l'espace, $(A)$
une sphère de centre $A$ et de rayon $r$, et $E$ l'ensemble des nombres
$R > 0$ tels qu'il existe une sphère $(H)$ de centre $H$ et de rayon $R$ par
rapport à laquelle les points $B$ et $C$ sont strictement extérieurs
(c'est à dire par exemple tels que $HB > R$), et les points de $(A)$
strictement intérieurs.

\medskip

\begin{enumerate}
\item Dans cette question, $B$ et $C$ sont alignés et strictement extérieurs à $(A)$.

Montrer que $E$ est non vide et majoré.
	Calculer le plus petit de ses majorants en fonction des données.
\item Déterminer une condition nécéssaire et suffisante pour $E$ soit non
	vide et majoré.
\item Calculer, lorsqu'il existe, le plus petit des majorants de $E$.
\end{enumerate}
\end{document}